\documentclass[11pt]{article}
\usepackage{mn}

\begin{document}
\zestaw{NUMER}{TEMAT}

\begin{teoria}[Jak czytać ten dokument]
Każde zadanie ma: przepisaną \textbf{treść} (szare pudełko), potrzebną
\textbf{teorię} (niebieskie), \textbf{intuicję} (żółte), \textbf{rozwiązanie
formalne} i~\textbf{odpowiedź} (zielone). Wyniki zostały sprawdzone
numerycznie skryptem \texttt{weryfikacja/NAZWA.py}. $\lg=\log_2$.
\end{teoria}

% ---------------------------------------------------------------------
\section{Zadanie 1. Krótki opis}
\begin{tresc}
Treść przepisana z zestawu.
\end{tresc}

\begin{teoria}[Potrzebne pojęcia]
\begin{definicja}...\end{definicja}
\begin{twierdzenie}...\end{twierdzenie}
\end{teoria}

\begin{intuicja}
Wyjaśnienie „na chłopski rozum”.
\end{intuicja}

\formalnie
Dowód / rachunki.

\begin{odpowiedz}
Wynik.
\end{odpowiedz}

\end{document}
